%% Generated by blueprint/scripts/render_registers.py from doc/DIVERGENCES.md. Do not edit.
\chapter{Divergences}\label{chap:divergences}

\section*{At a glance}
\begin{itemize}
\item \lead{Source} \code{doc/DIVERGENCES.md}, rendered by \code{blueprint/scripts/render\_registers.py}; the sections below are the register's own text.
\item \lead{Entries} 0.
\end{itemize}

This register lists every place where the verification code, or the statements it proves, diverge from the references of spec \S{}3.1--\S{}3.2 --- Letouzey, "A New Extraction for Coq" (\S{}3.1), and Sozeau et al., "Correct and Complete Type Checking and Certified Erasure for Coq, in Coq" (\S{}3.2), including its MetaRocq sources. The verification effort follows these references as closely as the differences between Lean's kernel/\texttt{lean4lean}'s model and Rocq/PCUIC/MetaRocq allow; every place it does not, for any reason, is recorded here.

Each entry has an id \texttt{DV-\textless{}n\textgreater{}} and exactly these fields:

\begin{itemize}
\item \textbf{Our artifact:} file and declaration on the \texttt{verification} branch that carries the divergence.
\item \textbf{Reference artifact:} the paper section/figure/theorem (\S{}3.1 or \S{}3.2) and the MetaRocq \texttt{file:identifier} it corresponds to.
\item \textbf{What differs:} the concrete difference between our artifact and the reference artifact.
\item \textbf{Why it is forced:} the concrete reason the divergence is necessary --- a difference between Lean's kernel and CIC, a gap in what \texttt{lean4lean} \texttt{master} provides, a scope boundary from \S{}2 of the spec, or similar. This field must name a real constraint; "it was simpler" or "we chose to" is not a forcing reason.
\item \textbf{What was considered instead:} the alternative(s) that would have kept the divergence from existing, and why each was rejected.
\end{itemize}

An entry without a forcing reason is a defect, not a divergence: it must be removed and the statement or definition restated to match the reference instead.

Entries are headed \texttt{\#\#\# DV-\textless{}n\textgreater{}}. They cite declarations by full name in backquotes (ours as \texttt{EraseProof.\textless{}name\textgreater{}}, lean4lean's as \texttt{Lean4Lean.\textless{}name\textgreater{}}), our files by their path in the repository (\texttt{proof/\allowbreak{}...}), lean4lean \texttt{master} sources by path and line (\texttt{Lean4Lean/\allowbreak{}...lean:\textless{}line\textgreater{}}), and MetaRocq sources by path relative to the MetaRocq repository and line. Check C9 of \texttt{proof/\allowbreak{}scripts/\allowbreak{}check.sh} fails if a cited \texttt{EraseProof.*} or \texttt{Lean4Lean.*} name is not a declaration, a cited file or line does not exist, an entry's fields differ from the list above, or an entry's artifact field cites no declaration of \texttt{EraseProof}.

The blueprint renders this register.

\section{\texorpdfstring{Entries}{Entries}}\label{reg:entries}

\subsection{\texorpdfstring{DV-1}{DV-1}}\label{reg:dv-1}

\begin{itemize}
\item \textbf{Our artifact:} \texttt{proof/\allowbreak{}EraseProof/\allowbreak{}Source/\allowbreak{}Eval.lean}, \texttt{EraseProof.SrcEval}: the source semantics evaluates Lean terms (\texttt{Expr}), the eraser's input, whose types are those of their images under \texttt{EraseProof.TrS} in lean4lean's syntax \texttt{VExpr}.
\item \textbf{Reference artifact:} PCUIC's weak call-by-value evaluation \texttt{eval} (\texttt{pcuic/\allowbreak{}theories/\allowbreak{}PCUICWcbvEval.v:231}; MetaCoq paper \S{}5.6) and the erasure relation \texttt{erases} (\texttt{erasure/\allowbreak{}theories/\allowbreak{}Extract.v:88}; \S{}7.3, Fig. 18), both on PCUIC terms, the syntax that PCUIC's typing judges; Letouzey \S{}3.1 (CIC terms).
\item \textbf{What differs:} the syntax that evaluates and the syntax that is typed are different: \texttt{Expr}, with \texttt{let} (\texttt{Expr.letE}) and metadata (\texttt{Expr.mdata}), for evaluation, and lean4lean's \texttt{VExpr} (\texttt{Lean4Lean/\allowbreak{}Theory/\allowbreak{}VExpr.lean:7-13}), which has neither, for typing; \texttt{EraseProof.TrS} relates them. In PCUIC both are the same terms. \texttt{EraseProof.SrcEval} evaluates a \texttt{let} as \texttt{eval\_\allowbreak{}zeta} does (\texttt{pcuic/\allowbreak{}theories/\allowbreak{}PCUICWcbvEval.v:241}): the value first, then the body with the value substituted; metadata is transparent (rule \texttt{EraseProof.SrcEval.mdata}).
\item \textbf{Why it is forced:} the eraser consumes \texttt{Expr}, and the \ensuremath{\lambda}\ensuremath{\Box} image of a \texttt{let} evaluates its value first (\texttt{eval\_\allowbreak{}zeta}, \texttt{erasure/\allowbreak{}theories/\allowbreak{}EWcbvEval.v:134}; rule \texttt{EraseProof.LBEval.zeta}). \texttt{VExpr} has no \texttt{let}: \texttt{EraseProof.TrS} translates \texttt{let x := v; b} to the translation of \texttt{b} in a context whose entry for \texttt{x} is the translation of \texttt{v} (rule \texttt{EraseProof.TrS.letE}, as rule \texttt{letE} of \texttt{Lean4Lean.TrExprS}, \texttt{Lean4Lean/\allowbreak{}Verify/\allowbreak{}Typing/\allowbreak{}Expr.lean:97-101}), that is, with the value substituted. A semantics on \texttt{VExpr} therefore never evaluates a \texttt{let} value, and relating it to \ensuremath{\lambda}\ensuremath{\Box} evaluation, which does, when the value diverges would need a normalization theorem, which lean4lean \texttt{master} does not have.
\item \textbf{What was considered instead:} a semantics on \texttt{VExpr}: with the recursive \texttt{unsafe def loop : A \ensuremath{\to} A := fun x =\textgreater{} loop x} and \texttt{a : A}, the term \texttt{let x := loop a; a} has the image \texttt{a}, a value, while its \ensuremath{\lambda}\ensuremath{\Box} image evaluates \texttt{loop a} first and diverges.
\end{itemize}

\subsection{\texorpdfstring{DV-2}{DV-2}}\label{reg:dv-2}

\begin{itemize}
\item \textbf{Our artifact:} \texttt{proof/\allowbreak{}EraseProof/\allowbreak{}Typing/\allowbreak{}Basic.lean}, \texttt{EraseProof.TrS}, the translation of Lean terms into lean4lean's model.
\item \textbf{Reference artifact:} MetaRocq's translation of Template terms into PCUIC terms, \texttt{trans} (\texttt{template-pcuic/\allowbreak{}theories/\allowbreak{}TemplateToPCUIC.v:63}; the MetaCoq paper's "certified type-preserving translation from Coq's syntax to PCUIC's syntax", \S{}9), which erasure (\S{}7) consumes; the translation \texttt{EraseProof.TrS} follows is lean4lean's \texttt{Lean4Lean.TrExprS} (\texttt{Lean4Lean/\allowbreak{}Verify/\allowbreak{}Typing/\allowbreak{}Expr.lean:75}).
\item \textbf{What differs:} \texttt{EraseProof.TrS} has the rules of \texttt{Lean4Lean.TrExprS} except \texttt{lit} and \texttt{proj}: literals (\texttt{Expr.lit}) and projections (\texttt{Expr.proj}) have no translation, so terms containing them are outside the fragment. Like \texttt{Lean4Lean.TrExprS} and unlike \texttt{trans}, it is a relation (its image is unique: \texttt{EraseProof.TrS.det}). Every derivation of \texttt{EraseProof.TrS} is one of \texttt{Lean4Lean.TrExprS} (test \texttt{EraseProof.Test.TrS.toTrExprS}, whose statement reaches the \texttt{sorry} of \texttt{Lean4Lean.TrProj}).
\item \textbf{Why it is forced:} the lemmas of \texttt{Lean4Lean.TrExprS} handle \texttt{proj} with \texttt{sorry} proofs (\texttt{Lean4Lean/\allowbreak{}Verify/\allowbreak{}Typing/\allowbreak{}Lemmas.lean:642,723,727,893,938,1240,1509}); reusing them puts projection sorries into a fragment without projections. The \texttt{lit} rule translates a literal through constructors of \texttt{Nat} and \texttt{String} (\texttt{Lean4Lean.VEnv.ContainsLits}, \texttt{Lean4Lean/\allowbreak{}Verify/\allowbreak{}Typing/\allowbreak{}Expr.lean:70}), which are inductive types, absent from the model (DV-6). Secondarily, projections are translated through \texttt{Lean4Lean.TrProj}, a \texttt{sorry} definition (\texttt{Lean4Lean/\allowbreak{}Verify/\allowbreak{}Typing/\allowbreak{}Expr.lean:68}), which every hypothesis stated with \texttt{Lean4Lean.TrExprS} would mention.
\item \textbf{What was considered instead:} \texttt{Lean4Lean.TrExprS} with a side condition excluding projections and literals: the hypotheses would still mention \texttt{Lean4Lean.TrProj}, and its lemmas would reach the \texttt{sorry} proofs of the \texttt{proj} cases.
\end{itemize}

\subsection{\texorpdfstring{DV-3}{DV-3}}\label{reg:dv-3}

\begin{itemize}
\item \textbf{Our artifact:} \texttt{proof/\allowbreak{}EraseProof/\allowbreak{}Env.lean}, \texttt{EraseProof.ProgEnv}, with \texttt{EraseProof.TrConst}, \texttt{EraseProof.TrDef} and \texttt{EraseProof.findDecl}: the relation between a program's declarations and their model in lean4lean.
\item \textbf{Reference artifact:} the well-formed global environment \texttt{wf\_\allowbreak{}ext \ensuremath{\Sigma}} (\texttt{pcuic/\allowbreak{}theories/\allowbreak{}PCUICTyping.v:507}; MetaCoq paper \S{}3.6), a hypothesis of \texttt{erases\_\allowbreak{}correct} (\texttt{erasure/\allowbreak{}theories/\allowbreak{}ErasureCorrectness.v:51}), whose declarations are a newest-first list searched by \texttt{lookup\_\allowbreak{}env} (\texttt{common/\allowbreak{}theories/\allowbreak{}Environment.v:483}); and lean4lean's relation between a kernel environment and its model, \texttt{Lean4Lean.TrEnv'} (\texttt{Lean4Lean/\allowbreak{}Verify/\allowbreak{}Environment/\allowbreak{}Basic.lean:128}), which \texttt{EraseProof.ProgEnv} restates.
\item \textbf{What differs:} the environment is a newest-first list of Lean \texttt{ConstantInfo}s, searched like \texttt{lookup\_\allowbreak{}env} by \texttt{EraseProof.findDecl}, instead of a Lean \texttt{Environment} or the constant map of \texttt{Lean4Lean.TrEnv'}. \texttt{EraseProof.ProgEnv} has the rules \texttt{axiom}, \texttt{defn}, \texttt{thm}, \texttt{opaque} and \texttt{mutualDef} (named \texttt{block}) of \texttt{Lean4Lean.TrEnv'} at safety \texttt{.unsafe}, with \texttt{EraseProof.TrS} in place of \texttt{Lean4Lean.TrExprS}; \texttt{EraseProof.TrConst} and \texttt{EraseProof.TrDef} restate \texttt{Lean4Lean.TrConstant} and \texttt{Lean4Lean.TrDefVal} without their safety conjunct, which holds for every declaration at \texttt{.unsafe}. It has no \texttt{ignore} rule (at \texttt{.unsafe} no declaration is ignored), no \texttt{quot} rule, no \texttt{induct} rule, and no freshness premise on a constant map (a name is fresh in the model because \texttt{Lean4Lean.VEnv.addConst} succeeds). It has one premise that \texttt{Lean4Lean.TrEnv'} lacks: the field \texttt{all} is what Lean's elaborator sets, the declaration's own name for a single definition, theorem or opaque constant, and the block's names in order for each member of a block. Every \texttt{EraseProof.ProgEnv} gives a \texttt{Lean4Lean.TrEnv'} at \texttt{.unsafe} of the same model, with a constant map that holds each of the program's declarations under its name (test \texttt{EraseProof.Test.ProgEnv.toTrEnv'}, whose statement reaches the \texttt{sorry} definitions \texttt{Lean4Lean.TrProj} and \texttt{Lean4Lean.VInductDecl.WF}). Compared with \texttt{wf\_\allowbreak{}ext \ensuremath{\Sigma}}, the environment has no inductive declarations, and its typing is lean4lean's (DV-5).
\item \textbf{Why it is forced:} a Lean \texttt{Environment} always contains the inductive types of \texttt{Init}, and \texttt{Lean4Lean.TrEnv'} at \texttt{.unsafe} relates no constant map that contains an inductive type (lean4lean's \texttt{TrEnv'.no\_\allowbreak{}inductInfo}, \texttt{Lean4Lean/\allowbreak{}Verify/\allowbreak{}Environment/\allowbreak{}Extension.lean:18}), because lean4lean \texttt{master} has no inductive types (DV-6); so a program is stated as an inductive-free list of declarations. The rule \texttt{quot} needs a constant named \texttt{Eq} (\texttt{Lean4Lean.VEnv.QuotReady}, \texttt{Lean4Lean/\allowbreak{}Theory/\allowbreak{}Quot.lean:13}), which is an inductive type in every Lean environment. \texttt{Lean4Lean.TrEnv'} is stated with \texttt{Lean4Lean.TrExprS}, whose projection rule and lemmas carry \texttt{sorry} (DV-2). The eraser reads \texttt{all} to lay out a block of mutual definitions (\texttt{visitMutual}, \texttt{LeanToLambdaBox/\allowbreak{}Erasure.lean}), and lean4lean's kernel does not check it (\texttt{addDefinition} to \texttt{addMutual}, \texttt{Lean4Lean/\allowbreak{}Environment.lean:36-118}), so the premise states what the elaborator guarantees.
\item \textbf{What was considered instead:} \texttt{Lean4Lean.TrEnv'} of a real Lean environment, which does not hold for any (\texttt{TrEnv'.no\_\allowbreak{}inductInfo}); \texttt{Lean4Lean.TrEnv'} of a kernel environment built from the program's declarations, whose hypotheses would mention \texttt{Lean4Lean.TrProj} through \texttt{Lean4Lean.TrExprS} (DV-2); no \texttt{all} premise, under which the eraser's reading of \texttt{all} is unconstrained.
\end{itemize}

\subsection{\texorpdfstring{DV-5}{DV-5}}\label{reg:dv-5}

\begin{itemize}
\item \textbf{Our artifact:} \texttt{proof/\allowbreak{}EraseProof/\allowbreak{}Typing/\allowbreak{}Basic.lean}: all typing in \texttt{EraseProof} is lean4lean's typing judgment \texttt{Lean4Lean.VEnv.IsDefEq} (through \texttt{Lean4Lean.VEnv.HasType}, \texttt{Lean4Lean.VEnv.IsType}, \texttt{Lean4Lean.VExpr.WF}), as in the premises of \texttt{EraseProof.TrS} and the conclusion of \texttt{EraseProof.TrS.wf}.
\item \textbf{Reference artifact:} PCUIC typing, MetaCoq paper \S{}3.3 (cumulativity) and \S{}3.5 (typing, Fig. 21), MetaRocq \texttt{pcuic/\allowbreak{}theories/\allowbreak{}PCUICTyping.v:198} \texttt{typing} with \texttt{type\_\allowbreak{}Cumul} (\texttt{:294}), on which erasure (\S{}7) is stated; Letouzey \S{}3.1 (CIC).
\item \textbf{What differs:} the judgment has definitional proof irrelevance (\texttt{proofIrrel}, \texttt{Lean4Lean/\allowbreak{}Theory/\allowbreak{}Typing/\allowbreak{}Basic.lean:51}) and \ensuremath{\eta} for functions (\texttt{eta}, \texttt{:48}; PCUIC has no \ensuremath{\eta}, MetaCoq paper \S{}2.4); it has no cumulativity (conversion only, \texttt{defeqDF}, \texttt{:44}); levels have the operator \texttt{imax} (\texttt{Lean4Lean/\allowbreak{}Theory/\allowbreak{}VLevel.lean:12}), in which the sort of a \ensuremath{\Pi}-type is stated (\texttt{forallEDF}, \texttt{Lean4Lean/\allowbreak{}Theory/\allowbreak{}Typing/\allowbreak{}Basic.lean:40}), and are compared up to evaluation (\texttt{Lean4Lean.VLevel.Equiv}, \texttt{Lean4Lean/\allowbreak{}Theory/\allowbreak{}VLevel.lean:58}, in \texttt{sortDF}, \texttt{Lean4Lean/\allowbreak{}Theory/\allowbreak{}Typing/\allowbreak{}Basic.lean:22}).
\item \textbf{Why it is forced:} the eraser's inputs are Lean terms, typed by Lean's kernel; the spec (\S{}2) requires lean4lean \texttt{master} as the model of that kernel, and it models Lean's theory, not PCUIC.
\item \textbf{What was considered instead:} none: PCUIC typing does not type Lean terms (it has cumulativity, and neither definitional proof irrelevance nor \texttt{imax}), so no judgment closer to the reference types the eraser's inputs.
\end{itemize}

\subsection{\texorpdfstring{DV-6}{DV-6}}\label{reg:dv-6}

\begin{itemize}
\item \textbf{Our artifact:} the fragment: the source terms \texttt{EraseProof.TrS} translates (\texttt{proof/\allowbreak{}EraseProof/\allowbreak{}Typing/\allowbreak{}Basic.lean}).
\item \textbf{Reference artifact:} MetaRocq's erasure relation \texttt{erases} in full (\texttt{erasure/\allowbreak{}theories/\allowbreak{}Extract.v:88}; MetaCoq paper \S{}7.3, Fig. 18), with its rules \texttt{erases\_\allowbreak{}tConstruct} (\texttt{:106}), \texttt{erases\_\allowbreak{}tCase} (\texttt{:109}), \texttt{erases\_\allowbreak{}tProj} (\texttt{:118}), \texttt{erases\_\allowbreak{}tFix} (\texttt{:122}), \texttt{erases\_\allowbreak{}tCoFix} (\texttt{:129}), \texttt{erases\_\allowbreak{}tPrim} (\texttt{:137}), and PCUIC's weak call-by-value evaluation \texttt{eval} (\texttt{pcuic/\allowbreak{}theories/\allowbreak{}PCUICWcbvEval.v:231}) with its \ensuremath{\iota}-rules; Letouzey Def. 3 (all cases of \ensuremath{\mathcal{E}}, fixpoints) and Def. 8 (\ensuremath{\iota}).
\item \textbf{What differs:} no constructors, case analysis, projections, source fixpoints, cofixpoints, primitive values or evars, and no \ensuremath{\iota}-reduction: \texttt{EraseProof.TrS} has no rule for \texttt{Expr.proj}, \texttt{Expr.lit} or \texttt{Expr.mvar}, and Lean's \texttt{Expr} has no node for constructors, case analysis, fixpoints or cofixpoints (in Lean these are constants of inductive declarations, which \texttt{EraseProof.TrS} translates only as constants of the model's environment).
\item \textbf{Why it is forced:} lean4lean \texttt{master} has no inductive types: \texttt{Lean4Lean.VInductDecl.WF} and \texttt{Lean4Lean.VEnv.addInduct} are \texttt{sorry} definitions (\texttt{Lean4Lean/\allowbreak{}Theory/\allowbreak{}Inductive.lean:5,7}), so the model has no constructors, recursors or \ensuremath{\iota}-reduction; projections are translated through the \texttt{sorry} definition \texttt{Lean4Lean.TrProj} (\texttt{Lean4Lean/\allowbreak{}Verify/\allowbreak{}Typing/\allowbreak{}Expr.lean:68}); literals are values of the inductive types \texttt{Nat} and \texttt{String} (DV-2); the kernel's terms have no metavariables.
\item \textbf{What was considered instead:} axiomatized inductive types (types, constructors and recursors as constants, \ensuremath{\iota} as definitional equalities) in an ordered environment: the rule \texttt{Lean4Lean.VEnv.Ordered.defeq} (\texttt{Lean4Lean/\allowbreak{}Theory/\allowbreak{}Typing/\allowbreak{}Lemmas.lean:258}) admits any well-typed definitional equality, such as \texttt{Sort 0 \ensuremath{\equiv} (Sort 0 \ensuremath{\to} Sort 0)}, so the injectivity of \ensuremath{\Pi}-types that the simulation proof needs fails there; master proves that injectivity (\texttt{Lean4Lean.VEnv.IsDefEqU.forallE\_\allowbreak{}inv}, \texttt{Lean4Lean/\allowbreak{}Theory/\allowbreak{}Typing/\allowbreak{}Injectivity.lean:23}) only under \texttt{Lean4Lean.VEnv.WF}, whose inductive part is the \texttt{sorry} definition above.
\end{itemize}

\subsection{\texorpdfstring{DV-7}{DV-7}}\label{reg:dv-7}

\begin{itemize}
\item \textbf{Our artifact:} \texttt{proof/\allowbreak{}EraseProof/\allowbreak{}Source/\allowbreak{}Eval.lean}: \texttt{EraseProof.RecursiveDecl} (with \texttt{EraseProof.OccursV}), the rules \texttt{EraseProof.SrcEval.fixAtom} and \texttt{EraseProof.SrcEval.fixApp} of \texttt{EraseProof.SrcEval}, the value shape \texttt{EraseProof.SrcValue.fixConst} and the head condition \texttt{EraseProof.BlocksCong}: recursion in the source semantics.
\item \textbf{Reference artifact:} PCUIC's fixpoints: the term \texttt{tFix}, a value (\texttt{atom}, \texttt{pcuic/\allowbreak{}theories/\allowbreak{}PCUICWcbvEval.v:51}), unfolded when applied by \texttt{eval\_\allowbreak{}fix} (\texttt{:273}) and excluded as a head of \texttt{eval\_\allowbreak{}app\_\allowbreak{}cong} (\texttt{:311}, \texttt{isFixApp}); their erasure \texttt{erases\_\allowbreak{}tFix} (\texttt{erasure/\allowbreak{}theories/\allowbreak{}Extract.v:122}); Letouzey Def. 3 (the fixpoint case of \ensuremath{\mathcal{E}}).
\item \textbf{What differs:} Lean has no fixpoint term; recursion lives in constants. A constant is recursive (\texttt{EraseProof.RecursiveDecl}) when its block has several members or its value mentions its own name in a value position (\texttt{EraseProof.OccursV}). \texttt{EraseProof.SrcEval} treats a recursive constant as a PCUIC \texttt{tFix} with \texttt{rarg = 0}: the constant evaluates to itself (\texttt{EraseProof.SrcEval.fixAtom}, \texttt{eval\_\allowbreak{}delta} to a \texttt{tFix} value), and applied to an argument it evaluates the argument, then its value, at the occurrence's levels, applied to the argument's value (\texttt{EraseProof.SrcEval.fixApp}, \texttt{eval\_\allowbreak{}fix} with \texttt{rarg = 0} and no earlier arguments). Recursive constants of every type are values, including those whose value is not a \ensuremath{\lambda}: with \texttt{unsafe def x : A \ensuremath{\to} A := x}, \texttt{x} is a value of \texttt{EraseProof.SrcEval}, as its image, a \texttt{tFix}, is a value of \ensuremath{\lambda}\ensuremath{\Box}, while Lean's own evaluation of \texttt{x} (\texttt{lean -{}-run}) does not terminate. Non-recursive constants unfold when evaluated (\texttt{EraseProof.SrcEval.delta}, \texttt{eval\_\allowbreak{}delta}). Which constants are recursive is a syntactic property of the declaration.
\item \textbf{Why it is forced:} Lean's \texttt{Expr} has no fixpoint node: a recursive Lean definition is a constant whose value mentions itself, a member of a block (rule \texttt{block} of \texttt{EraseProof.ProgEnv}, whose values are translated in the model that contains the block). The eraser compiles a declaration to a member of a \ensuremath{\lambda}\ensuremath{\Box} \texttt{tFix} exactly when it is not a single declaration whose value does not mention its own name (\texttt{name\_\allowbreak{}occurs} in \texttt{visitMutual}, \texttt{LeanToLambdaBox/\allowbreak{}Erasure.lean}), and \ensuremath{\lambda}\ensuremath{\Box} has fixpoint values only for these blocks. For the source semantics to be simulated, the constants it treats as fixpoints must be exactly these, so the split is the eraser's test, which \texttt{EraseProof.RecursiveDecl} states on the declaration.
\item \textbf{What was considered instead:} unfolding recursive constants when evaluated, as the others: it matches \ensuremath{\lambda}\ensuremath{\Box}'s \texttt{tFix} only when each member's value is a \ensuremath{\lambda}, so that one unfolding reaches a value; it would need a shipping change that \ensuremath{\eta}-expands members whose value is not a \ensuremath{\lambda}, which the theorem does not need (spec \S{}5.1).
\end{itemize}

\subsection{\texorpdfstring{DV-10}{DV-10}}\label{reg:dv-10}

\begin{itemize}
\item \textbf{Our artifact:} \texttt{proof/\allowbreak{}EraseProof/\allowbreak{}Source/\allowbreak{}EvalEnv.lean}, \texttt{EraseProof.EvalEnv.unfold?}, the constants the source semantics unfolds: it returns the value of a definition or theorem and nothing for an opaque constant.
\item \textbf{Reference artifact:} the \ensuremath{\delta} rule of PCUIC's weak call-by-value evaluation, \texttt{eval\_\allowbreak{}delta} (\texttt{pcuic/\allowbreak{}theories/\allowbreak{}PCUICWcbvEval.v:247}; MetaCoq paper \S{}5.6), which unfolds every constant whose declaration has a body (\texttt{cst\_\allowbreak{}body decl = Some body}).
\item \textbf{What differs:} a Lean \texttt{opaque} declaration (\texttt{ConstantInfo.opaqueInfo}) has a value, but \texttt{EraseProof.EvalEnv.unfold?} does not return it, so evaluation does not unfold the constant. The shipping eraser emits the opaque's value as the body of its \ensuremath{\lambda}\ensuremath{\Box} constant (\texttt{ci.value? (allowOpaque := true)} in \texttt{LeanToLambdaBox/\allowbreak{}Erasure.lean}), so \ensuremath{\lambda}\ensuremath{\Box} evaluation unfolds what the source semantics leaves stuck.
\item \textbf{Why it is forced:} lean4lean's model gives an opaque constant no defining equation (rule \texttt{opaque} of \texttt{Lean4Lean.TrEnv'}, \texttt{Lean4Lean/\allowbreak{}Verify/\allowbreak{}Environment/\allowbreak{}Basic.lean:164-169}, adds the constant only), and \texttt{EraseProof.ProgEnv} follows it; so unfolding an opaque is not a definitional equality of the model. The proof relates each evaluation of the source to a typed definitional equality of the model (\texttt{EraseProof.SrcEval.defeq}, subject reduction); for a \ensuremath{\delta} step, \texttt{EraseProof.ProgEnv.unfold} gives it for definitions (their defining equation) and theorems (their type is a proposition, so proof irrelevance equates them with their value), and nothing gives it for opaques, whose type need not be a proposition.
\item \textbf{What was considered instead:} unfolding opaques as \texttt{eval\_\allowbreak{}delta} does: evaluation steps would no longer be definitional equalities of the model, and type preservation fails for terms whose type depends on an opaque's value (with \texttt{opaque c : T := t} and an axiom \texttt{f : (x : T) \ensuremath{\to} F x}, \texttt{f c : F c} evaluates to \texttt{f t : F t}, and the model does not equate \texttt{F c} with \texttt{F t}).
\end{itemize}

\subsection{\texorpdfstring{DV-11}{DV-11}}\label{reg:dv-11}

\begin{itemize}
\item \textbf{Our artifact:} \texttt{proof/\allowbreak{}EraseProof/\allowbreak{}Atoms.lean}, \texttt{EraseProof.EvalEnv.isAtom}, with \texttt{EraseProof.arityShape}, \texttt{EraseProof.EvidentZero}, \texttt{EraseProof.DeltaFree} and \texttt{EraseProof.EvidentProp}; \texttt{proof/\allowbreak{}EraseProof/\allowbreak{}Source/\allowbreak{}Eval.lean}, the rule \texttt{EraseProof.SrcEval.constAtom} and the values \texttt{EraseProof.AtomSpine} (in \texttt{EraseProof.SrcValue.spine}): the constants without a \ensuremath{\delta} rule that are values of the source semantics.
\item \textbf{Reference artifact:} PCUIC's atoms \texttt{atom} (\texttt{pcuic/\allowbreak{}theories/\allowbreak{}PCUICWcbvEval.v:51}), which contain the inductive types \texttt{tInd} and the constructors \texttt{tConstruct}; \texttt{eval\_\allowbreak{}atom} (\texttt{:331}) evaluates them to themselves, \texttt{eval\_\allowbreak{}app\_\allowbreak{}cong} (\texttt{:311}) evaluates their applications, and \texttt{value} (\texttt{:500}) lists the results; a constant without a body has no rule (\texttt{eval\_\allowbreak{}delta}, \texttt{:247}, needs \texttt{cst\_\allowbreak{}body decl = Some body}). On the erasure side, \texttt{erases} (\texttt{erasure/\allowbreak{}theories/\allowbreak{}Extract.v:88}) has no rule for \texttt{tInd}, which erases only to \texttt{\ensuremath{\Box}} (\texttt{erases\_\allowbreak{}box}, \texttt{:140}), and \texttt{erases\_\allowbreak{}tConstruct} (\texttt{:106}) requires the inductive not to be propositional (\texttt{isPropositional}, \texttt{pcuic/\allowbreak{}theories/\allowbreak{}PCUICFirstorder.v:109}, which reads the inductive's declared arity syntactically: \texttt{isPropositionalArity}, \texttt{:103}, with \texttt{destArity}, \texttt{pcuic/\allowbreak{}theories/\allowbreak{}PCUICAst.v:486}).
\item \textbf{What differs:} a constant without a \ensuremath{\delta} rule in the source semantics is a value when its declared type is evidently an arity (\texttt{EraseProof.arityShape}: syntactically \texttt{\ensuremath{\Pi} x\ensuremath{_1}\ldots{}x\ensuremath{_n}, Sort u}) or evidently a proposition (\texttt{EraseProof.EvidentProp}): a \ensuremath{\Pi}-telescope over an application \texttt{h.\{vs\} a\ensuremath{_1}\ldots{}a\ensuremath{_n}} whose head \texttt{h} is an axiom or an opaque constant (\texttt{EraseProof.DeltaFree}) with declared type a syntactic arity \texttt{\ensuremath{\Pi} y\ensuremath{_1}\ldots{}y\ensuremath{_n}, Sort u} of exactly \texttt{n} binders, and whose sort at that occurrence, \texttt{u} with \texttt{h}'s level parameters instantiated by \texttt{vs} (\texttt{Erasure.Pure.instLevel}), is evidently zero (\texttt{EraseProof.EvidentZero}). Such a constant evaluates to itself (\texttt{EraseProof.SrcEval.constAtom}), and the applications it heads are values (\texttt{EraseProof.AtomSpine}). Atomhood is a property of the constant, decided on declared types, not on the levels of an occurrence: a proof \texttt{hq.\{v\} : P.\{v\}} (with \texttt{axiom P.\{v\} : Sort v}) whose propositionality depends on its own level parameter is not an atom, not even at \texttt{hq.\{0\}}. Every other constant without a \ensuremath{\delta} rule is stuck.
\item \textbf{Why it is forced:} the fragment has no inductive types (DV-6), so its base types and the canonical proofs of atomic propositions are axioms; without a rule that makes them values, no program that instantiates a polymorphic function at a base type or passes such a proof by value evaluates, and the correctness theorem says nothing about it. The class is syntactic ("evident") because the eraser must box each atom from the declared types alone, without a typing hypothesis: a semantic class (every erasable constant without a \ensuremath{\delta} rule) would need the eraser's oracle \texttt{Erasure.Pure.isErasable} to box all of them, which needs canonicity (no neutral term is definitionally equal to a sort), and lean4lean \texttt{master} does not prove it. The head is \texttt{EraseProof.DeltaFree} because the head of a propositional constructor's type is an inductive type, which never \ensuremath{\delta}-reduces: a proof whose proposition is headed by a definition that unfolds to a \ensuremath{\Pi} can be applied, and the oracle then answers from the arguments' types (with \texttt{axiom hq : Q}, \texttt{def Q : Prop := \ensuremath{\forall} P : Prop, P \ensuremath{\to} P}, \texttt{axiom A : Type} and \texttt{axiom a : A}, it keeps the ill-typed spine \texttt{hq A a}). The head's sort is read at the occurrence because a Lean level can be \texttt{0}, whereas \texttt{isPropositional} reads an inductive's declared sort, which no Rocq universe instance turns into \texttt{Prop}. Atomhood does not depend on the occurrence because the eraser erases a universe-polymorphic body once, at its level parameters, where the oracle keeps the level-dependent proof \texttt{hq.\{v\}} (it boxes \texttt{hq.\{0\}}): were \texttt{hq.\{0\}} a value, a body passing \texttt{hq.\{v\}}, evaluated at level \texttt{0}, would reach a value in the source while its \ensuremath{\lambda}\ensuremath{\Box} image is stuck on the axiom \texttt{hq}, and erasure would not commute with level instantiation, as \texttt{erases\_\allowbreak{}subst\_\allowbreak{}instance\_\allowbreak{}decl} (\texttt{erasure/\allowbreak{}theories/\allowbreak{}ErasureProperties.v:412}) states it does.
\item \textbf{What was considered instead:} every constant without a \ensuremath{\delta} rule stuck (the theorem then says nothing about the programs above); a semantic class (needs canonicity); heads that are definitions (the oracle keeps proofs with such heads, above); heads that \ensuremath{\delta}-reduce to an \texttt{EraseProof.DeltaFree} head (further from the reference, and the eraser's boxing of the class would have to follow the oracle's \ensuremath{\delta} steps through definitions); the head's declared sort instead of its sort at the occurrence (leaves proofs of \texttt{P.\{0\}}, with \texttt{axiom P.\{v\} : Sort v}, stuck with no forcing reason); atomhood decided at each occurrence (makes the theorem false, above); the head's declared type read up to the kernel's \ensuremath{\delta} (the class would no longer be read syntactically, as \texttt{isPropositionalArity} reads the declared arity, and the eraser's boxing of the class would have to follow the oracle's \ensuremath{\delta} steps through definitions).
\end{itemize}

\subsection{\texorpdfstring{DV-12}{DV-12}}\label{reg:dv-12}

\begin{itemize}
\item \textbf{Our artifact:} \texttt{proof/\allowbreak{}EraseProof/\allowbreak{}Source/\allowbreak{}EvalEnv.lean}, \texttt{EraseProof.EvalEnv} with its field \texttt{axiomatized}, and \texttt{EraseProof.EvalEnv.unfold?}: the evaluation environment carries the constants the configuration remaps to foreign code, and evaluation does not unfold them.
\item \textbf{Reference artifact:} the environment of PCUIC's weak call-by-value evaluation \texttt{eval} (\texttt{pcuic/\allowbreak{}theories/\allowbreak{}PCUICWcbvEval.v:231}), whose rule \texttt{eval\_\allowbreak{}delta} (\texttt{:247}) unfolds a constant exactly when its declaration has a body; and \texttt{erases\_\allowbreak{}constant\_\allowbreak{}body} (\texttt{erasure/\allowbreak{}theories/\allowbreak{}Extract.v:264}), which erases a constant with a body to one with a body and a constant without a body to one without.
\item \textbf{What differs:} a definition or theorem for which \texttt{axiomatized} holds has a Lean value but no \ensuremath{\delta} rule in the source semantics: it is stuck, like the \ensuremath{\lambda}\ensuremath{\Box} axiom the eraser emits for it. Which constants are remapped is an input of the source semantics, besides the declarations.
\item \textbf{Why it is forced:} under the configuration \texttt{extern := .preferAxiom} (the default of \texttt{ErasureConfig}, \texttt{LeanToLambdaBox/\allowbreak{}Erasure.lean}), the shipping eraser emits a declaration tagged \texttt{@[extern]} as a \ensuremath{\lambda}\ensuremath{\Box} axiom, although it has a Lean value, so that it is linked with a foreign implementation (Dima \S{}4.2). Neither lean4lean's model nor \ensuremath{\lambda}\ensuremath{\Box} evaluation models that implementation; the source semantics can only leave the constant stuck, as \ensuremath{\lambda}\ensuremath{\Box} evaluation leaves the axiom.
\item \textbf{What was considered instead:} a scope condition excluding programs that mention a remapped constant: it also excludes programs that never evaluate the constant.
\end{itemize}

\subsection{\texorpdfstring{DV-16}{DV-16}}\label{reg:dv-16}

\begin{itemize}
\item \textbf{Our artifact:} \texttt{proof/\allowbreak{}EraseProof/\allowbreak{}Target.lean}, \texttt{EraseProof.LBEval}, \ensuremath{\lambda}\ensuremath{\Box} weak call-by-value evaluation, with its atoms \texttt{EraseProof.lbAtom}, stated on the shipping AST \texttt{LBTerm} (\texttt{LeanToLambdaBox/\allowbreak{}Basic.lean}); with it the functions \texttt{EraseProof.csubst}, \texttt{EraseProof.closedn}, \texttt{EraseProof.hasFVar} and the environment condition \texttt{EraseProof.LenvClosed}.
\item \textbf{Reference artifact:} MetaRocq's \ensuremath{\lambda}\ensuremath{\Box} evaluation \texttt{eval} (\texttt{erasure/\allowbreak{}theories/\allowbreak{}EWcbvEval.v:119}) on the \ensuremath{\lambda}\ensuremath{\Box} terms \texttt{term} (\texttt{erasure/\allowbreak{}theories/\allowbreak{}EAst.v:29}), with \texttt{atom} (\texttt{erasure/\allowbreak{}theories/\allowbreak{}EWcbvEval.v:36}), \texttt{csubst} (\texttt{erasure/\allowbreak{}theories/\allowbreak{}ECSubst.v:14}), \texttt{closedn} (\texttt{erasure/\allowbreak{}theories/\allowbreak{}ELiftSubst.v:90}) and \texttt{closed\_\allowbreak{}env} (\texttt{erasure/\allowbreak{}theories/\allowbreak{}EGlobalEnv.v:181}); MetaCoq paper \S{}7.1 (Fig. 16 and the amended evaluation rules, p. 8:60).
\item \textbf{What differs:} \texttt{LBTerm} has no \texttt{tVar}, \texttt{tEvar}, \texttt{tCoFix}, \texttt{tLazy} or \texttt{tForce}, and its only primitive values are 63-bit integers. So \texttt{EraseProof.LBEval} has no rules \texttt{eval\_\allowbreak{}cofix\_\allowbreak{}case} (\texttt{erasure/\allowbreak{}theories/\allowbreak{}EWcbvEval.v:198}), \texttt{eval\_\allowbreak{}cofix\_\allowbreak{}proj} (\texttt{:205}) or \texttt{eval\_\allowbreak{}force} (\texttt{:279}), its rule \texttt{prim} (\texttt{eval\_\allowbreak{}prim}, \texttt{:275}) evaluates an integer to itself, and \texttt{EraseProof.lbAtom} has no case for \texttt{tCoFix} or \texttt{tLazy}. \texttt{LBTerm} has a constructor that \texttt{term} lacks, \texttt{fvar}, a free variable named by a Lean \texttt{FVarId}: \texttt{EraseProof.LBEval} has no rule for it, \texttt{EraseProof.csubst} leaves it unchanged and \texttt{EraseProof.closedn} counts it as closed, as MetaRocq's \texttt{csubst} and \texttt{closedn} treat \texttt{tVar}; \texttt{EraseProof.hasFVar} tests whether it occurs, and has no MetaRocq counterpart; \texttt{EraseProof.LenvClosed} requires, besides the closedness of \texttt{closed\_\allowbreak{}env}, that no body stored in the \ensuremath{\lambda}\ensuremath{\Box} environment contains a free variable. The rules for constructors in block form, \texttt{eval\_\allowbreak{}iota\_\allowbreak{}block} (\texttt{:151}), \texttt{eval\_\allowbreak{}proj\_\allowbreak{}block} (\texttt{:228}) and \texttt{eval\_\allowbreak{}construct\_\allowbreak{}block} (\texttt{:254}), are absent; each requires the flag \texttt{with\_\allowbreak{}constructor\_\allowbreak{}as\_\allowbreak{}block} to be true, which it is not in \texttt{default\_\allowbreak{}wcbv\_\allowbreak{}flags} (\texttt{:69}), the flags of \texttt{erases\_\allowbreak{}correct} (\texttt{erasure/\allowbreak{}theories/\allowbreak{}ErasureCorrectness.v:51}). At flags where that flag is false, \texttt{EraseProof.LBEval} has exactly the rules of \texttt{eval} whose term constructors \texttt{LBTerm} has.
\item \textbf{Why it is forced:} the theorem is about the shipping eraser (spec \S{}2), whose output is an \texttt{LBTerm}, the \ensuremath{\lambda}\ensuremath{\Box} that peregrine reads; evaluation is stated on that type. The terms without a constructor in \texttt{LBTerm} cannot occur in the eraser's output, so their rules have nothing to apply to. The constructor \texttt{fvar} is part of the shipping AST because the traversal is locally nameless (DV-13): it opens binders with free variables and closes them with \texttt{abstract} and \texttt{toBvar} (\texttt{LeanToLambdaBox/\allowbreak{}Basic.lean}) before a body is stored, which \texttt{EraseProof.LenvClosed} records. MetaRocq's erasure works on de Bruijn indices and maps \texttt{tVar} only to \texttt{tVar} (\texttt{erasure/\allowbreak{}theories/\allowbreak{}ErasureFunction.v:996}), which typed terms do not contain (PCUIC's typing has no rule for it), so \texttt{closed\_\allowbreak{}env} needs no such clause. The block rules cannot fire at the flags of \texttt{erases\_\allowbreak{}correct}, at which the simulation evaluates; results stated for every flag, such as \texttt{EraseProof.LBEval.closed}, are about the relation without them.
\item \textbf{What was considered instead:} extending \texttt{LBTerm} with the missing constructors, a shipping change that no output of the eraser needs (spec \S{}5.1 allows only strictly necessary shipping changes); stating evaluation on a separate transcription of \texttt{term} reached through a translation from \texttt{LBTerm}, which every statement would pass through and which would still need an image for \texttt{fvar}; including the block rules, which the statements never use at their flags.
\end{itemize}

\subsection{\texorpdfstring{DV-21}{DV-21}}\label{reg:dv-21}

\begin{itemize}
\item \textbf{Our artifact:} \texttt{proof/\allowbreak{}EraseProof/\allowbreak{}Env.lean}, \texttt{EraseProof.ProgEnv}: its rule \texttt{axiom} admits axioms, declarations without a value, in the program's environment; and \texttt{proof/\allowbreak{}EraseProof/\allowbreak{}Source/\allowbreak{}Eval.lean}, \texttt{EraseProof.SrcEval}, which evaluates programs that use them.
\item \textbf{Reference artifact:} Letouzey \S{}3.4 (p. 10), "from now to the end of this paper we will only consider contexts with no assumptions", the hypothesis of Theorems 12, 13 and 15. MetaRocq's \texttt{erases\_\allowbreak{}correct} (\texttt{erasure/\allowbreak{}theories/\allowbreak{}ErasureCorrectness.v:51}) has no such hypothesis: \texttt{wf\_\allowbreak{}ext \ensuremath{\Sigma}} admits constants without a body, and \texttt{axiom\_\allowbreak{}free} (\texttt{erasure/\allowbreak{}theories/\allowbreak{}Extract.v:381}) is a hypothesis only of the first-order results, such as \texttt{erase\_\allowbreak{}correct\_\allowbreak{}firstorder} (\texttt{erasure/\allowbreak{}theories/\allowbreak{}ErasureFunctionProperties.v:2310}).
\item \textbf{What differs:} programs may depend on axioms: \texttt{EraseProof.ProgEnv} relates lists containing Lean \texttt{axiom} declarations to models in which they are constants without a defining equation. In evaluation an axiom has no \ensuremath{\delta} rule: an axiom whose declared type is evidently an arity or a proposition (\texttt{EraseProof.EvalEnv.isAtom}, DV-11) is a value (\texttt{EraseProof.SrcEval.constAtom}), and any other axiom is stuck (no rule of \texttt{EraseProof.SrcEval} evaluates it), as a constant without a body is in PCUIC's \texttt{eval} (\texttt{pcuic/\allowbreak{}theories/\allowbreak{}PCUICWcbvEval.v:231}). This follows MetaRocq and diverges from Letouzey.
\item \textbf{Why it is forced:} the spec (\S{}2) puts in scope every input whose verification needs nothing beyond lean4lean \texttt{master}, and \texttt{master} models axioms (rule \texttt{axiom} of \texttt{Lean4Lean.TrEnv'}, \texttt{Lean4Lean/\allowbreak{}Verify/\allowbreak{}Environment/\allowbreak{}Basic.lean:134}). Letouzey needs the hypothesis for canonicity (a closed term of an inductive type reduces to a constructor), in the \ensuremath{\iota} cases of the proof of Theorem 12 (Appendix A, cases 1 and 2) and in Theorem 15; the fragment has no inductive types and no \ensuremath{\iota}-reduction (DV-6).
\item \textbf{What was considered instead:} excluding axioms from \texttt{EraseProof.ProgEnv}, as Letouzey does: it narrows the scope the spec fixes without a gap in \texttt{master} that forces it.
\end{itemize}

\subsection{\texorpdfstring{DV-22}{DV-22}}\label{reg:dv-22}

\begin{itemize}
\item \textbf{Our artifact:} \texttt{proof/\allowbreak{}EraseProof/\allowbreak{}Source/\allowbreak{}Eval.lean}, \texttt{EraseProof.SrcEval}: the source semantics is big-step, call-by-value evaluation.
\item \textbf{Reference artifact:} Letouzey \S{}3.4, Definition 9 (weak reductions: the \ensuremath{\beta}, \ensuremath{\iota}, \ensuremath{\delta}, \ensuremath{\zeta} and \ensuremath{\Box} steps, closed under application, on either side, and under case analysis, but not under \ensuremath{\lambda}, so in any order), for which Theorem 13 (backward simulation) is stated step by step. MetaRocq's \texttt{erases\_\allowbreak{}correct} (\texttt{erasure/\allowbreak{}theories/\allowbreak{}ErasureCorrectness.v:51}) is stated for PCUIC's big-step weak call-by-value evaluation \texttt{eval} (\texttt{pcuic/\allowbreak{}theories/\allowbreak{}PCUICWcbvEval.v:231}; MetaCoq paper \S{}5.6), which \texttt{EraseProof.SrcEval} follows.
\item \textbf{What differs:} \texttt{EraseProof.SrcEval} fixes one order, call-by-value: an application evaluates its function and its argument before the \ensuremath{\beta}-step (\texttt{EraseProof.SrcEval.beta}), and a \texttt{let} its value before its body (\texttt{EraseProof.SrcEval.zeta}). Reductions in any other order, which Definition 9 admits, are not steps of the source semantics, which is big-step: it relates a term to its value, not to its one-step reducts.
\item \textbf{Why it is forced:} the target is \ensuremath{\lambda}\ensuremath{\Box}'s call-by-value evaluation (\texttt{EraseProof.LBEval}, MetaRocq's \texttt{eval}, \texttt{erasure/\allowbreak{}theories/\allowbreak{}EWcbvEval.v:119}, which peregrine implements), and a value reached in another order need not be reached by it: with the recursive \texttt{unsafe def loop : A \ensuremath{\to} A := fun x =\textgreater{} loop x} and \texttt{a : A}, \texttt{(fun \_\allowbreak{} =\textgreater{} a) (loop a)} reduces to \texttt{a} by one \ensuremath{\beta}-step, while \ensuremath{\lambda}\ensuremath{\Box} evaluation of its image evaluates \texttt{loop a} first and diverges. So the theorem, which is about the \ensuremath{\lambda}\ensuremath{\Box} evaluator, needs a source semantics that evaluates in the evaluator's order; MetaRocq's \texttt{erases\_\allowbreak{}correct} is stated in the same way.
\item \textbf{What was considered instead:} weak reduction in any order, as Definition 9, with a step-by-step simulation, as Theorem 13: it relates single steps and says nothing about the value the call-by-value evaluator reaches, and a value reached in another order need not be reached by it (above).
\end{itemize}

